\chapter{MeshGraphNet-часть для прогнозирования термоупругих полей на конечно-элементной сетке}

\section{Назначение MeshGraphNet-части в общей дипломной работе}

Данная глава описывает индивидуальную часть дипломной работы, связанную с применением графовой нейронной сети MeshGraphNet для прогнозирования распространения термоупругих полей в геологических средах. В общей структуре работы COMSOL Multiphysics используется как инженерный инструмент для построения физически корректной численной модели, задания материалов, граничных условий, начальных условий, конечно-элементной сетки и получения расчётных данных. MeshGraphNet-часть не заменяет COMSOL-моделирование на этапе построения эталонных симуляций, а использует экспортированные из COMSOL данные как обучающий материал для нейросетевой модели.

Основная задача MeshGraphNet-части заключается в следующем:
\begin{itemize}
    \item загрузить конечно-элементную сетку, экспортированную из COMSOL;
    \item загрузить физические поля, полученные в результате расчёта термоупругой задачи;
    \item преобразовать конечно-элементную сетку в графовую структуру данных;
    \item сформировать узловые и рёберные признаки для графовой нейронной сети;
    \item построить обучающий датасет из последовательности временных состояний;
    \item обучить MeshGraphNet прогнозировать развитие физических полей во времени;
    \item выполнить одношаговое и многошаговое прогнозирование на основе обученной модели;
    \item оценить качество прогноза относительно расчётных данных COMSOL.
\end{itemize}

В рассматриваемой постановке COMSOL является источником двух основных типов информации. Первый тип --- это геометрическая и топологическая информация о конечно-элементной сетке. Она хранится в файле \texttt{sandstone.mphtxt}. Второй тип --- это расчётные физические поля, сохранённые в CSV-файлах: температура, перемещения, скорости, свойства материалов, деформации и напряжения. MeshGraphNet использует эти данные для обучения оператора перехода от состояния системы в момент времени $t$ к состоянию системы в следующий момент времени $t+1$.

В общем виде численный решатель COMSOL задаёт переход между состояниями физической системы:
\begin{equation}
    X^{t+1}=\Phi_{\mathrm{COMSOL}}\left(X^t,U^t,M,\mathcal{M}\right),
\end{equation}
где $X^t$ --- состояние физических полей в момент времени $t$, $U^t$ --- внешние воздействия и граничные условия, $M$ --- параметры материалов, $\mathcal{M}$ --- конечно-элементная сетка. MeshGraphNet обучается аппроксимировать этот переход на графе:
\begin{equation}
    \hat{X}^{t+1}=F_{\theta}\left(G^t\right),
\end{equation}
где $G^t$ --- графовое представление состояния конечно-элементной модели в момент времени $t$, $F_{\theta}$ --- параметризованная графовая нейронная сеть, $\theta$ --- обучаемые параметры модели.

Таким образом, MeshGraphNet-часть является отдельным программно-исследовательским модулем дипломной работы. Она отвечает не за создание физической модели в COMSOL, а за машинное обучение на данных этой модели: подготовку графового датасета, обучение нейросетевого оператора и прогнозирование термоупругой динамики на неструктурированной сетке.

\section{Причина выбора графового представления}

Термоупругая задача в COMSOL решается методом конечных элементов. При таком подходе геометрическая область разбивается на конечные элементы, например тетраэдры в трёхмерном случае. Узлы элементов образуют дискретное представление геологической среды и нагретого металлического стержня, а физические величины вычисляются в узлах, элементах или интерполируются между ними.

Для данной задачи графовое представление является естественным по нескольким причинам. Во-первых, конечно-элементная сетка обычно является неструктурированной. Это означает, что узлы не расположены в виде регулярной декартовой решётки, как в изображениях или классических трёхмерных массивах. Расстояния между соседними узлами могут различаться, а число соседей у разных узлов может быть неодинаковым. Во-вторых, физическое взаимодействие в конечно-элементной модели определяется не только координатами точек, но и связностью сетки: какие узлы входят в один элемент, какие узлы являются соседними, через какие локальные области передаётся тепло и механическое воздействие. В-третьих, при переводе такой сетки в обычную таблицу теряется важная топологическая информация, потому что строки таблицы сами по себе не описывают соседство узлов.

Граф позволяет явно сохранить эту структуру. Узлы графа соответствуют узлам конечно-элементной сетки, а рёбра графа соответствуют локальным связям между узлами, возникающим из конечных элементов. Тогда распространение тепла, перемещений, деформаций и напряжений можно приближённо описывать через обмен информацией между соседними узлами. Именно такую логику использует MeshGraphNet: модель выполняет несколько итераций передачи сообщений по рёбрам графа, благодаря чему каждый узел получает информацию от своей локальной окрестности.

Преимущество графового представления состоит в том, что оно не требует приведения данных к регулярной сетке. Это особенно важно для конечно-элементных симуляций, где плотность сетки может быть повышена в областях с сильными градиентами температуры, напряжений или деформаций. В рассматриваемой задаче такие области могут возникать около нагретого металлического стержня и вблизи зон интенсивного термоупругого отклика геологической среды.

\section{Исходные данные для MeshGraphNet}

Исходные данные для MeshGraphNet формируются на основе результатов расчёта в COMSOL. После выполнения симуляции экспортируются сетка и таблицы физических величин. В данной работе используются файлы, перечисленные в таблице~\ref{tab:mgn_input_files}.

\begin{table}[h]
\centering
\caption{Исходные файлы MeshGraphNet-части}
\label{tab:mgn_input_files}
\begin{tabular}{|p{0.28\textwidth}|p{0.62\textwidth}|}
\hline
\textbf{Файл} & \textbf{Назначение} \\
\hline
\texttt{sandstone.mphtxt} & Файл конечно-элементной сетки COMSOL. Содержит координаты узлов, описание конечных элементов и топологию сетки, на основе которой строится граф. \\
\hline
\texttt{data\_temperature.csv} & Таблица температурного поля. Содержит координаты узлов или точек экспорта и значение температуры $T$ по временным шагам. \\
\hline
\texttt{data\_displacement.csv} & Таблица механического состояния. Содержит перемещения $u$, $v$, $w$ и, при наличии в экспорте, скорости $u_t$, $v_t$, $w_t$. \\
\hline
\texttt{data\_materials.csv} & Таблица свойств материалов. Содержит теплопроводность, плотность, теплоёмкость, модуль Юнга, коэффициент Пуассона, коэффициент температурного расширения и другие параметры, необходимые для переноса модели между материалами. \\
\hline
\texttt{data\_strain.csv} & Таблица деформаций. Содержит компоненты деформации, например $\varepsilon_x$, $\varepsilon_y$, $\varepsilon_z$. \\
\hline
\texttt{data\_stress\_1.csv} & Таблица компонент напряжений. Может содержать часть компонент тензора напряжений, экспортированных из COMSOL. \\
\hline
\texttt{data\_stress\_2.csv} & Дополнительная таблица компонент напряжений. Используется, если компоненты напряжений были сохранены несколькими отдельными файлами. \\
\hline
\texttt{data\_stress\_3.csv} & Дополнительная таблица компонент напряжений или связанных механических величин. Объединяется с остальными таблицами по координатам, идентификаторам узлов и времени. \\
\hline
\end{tabular}
\end{table}

Набор используемых столбцов зависит от того, какие величины были экспортированы из COMSOL. В MeshGraphNet-части используется следующая логика: координаты и статические свойства материалов используются как входные признаки, а физические поля, изменяющиеся во времени, используются как входные признаки и как целевые переменные для прогноза следующего временного шага.

Основные поля, используемые при построении датасета, приведены в таблице~\ref{tab:mgn_fields}.

\begin{table}[h]
\centering
\caption{Физические поля, используемые в MeshGraphNet-датасете}
\label{tab:mgn_fields}
\begin{tabular}{|p{0.28\textwidth}|p{0.62\textwidth}|}
\hline
\textbf{Поле} & \textbf{Физический смысл} \\
\hline
\texttt{x}, \texttt{y}, \texttt{z} & Пространственные координаты узла конечно-элементной сетки. \\
\hline
\texttt{T} & Температура в узле или точке экспорта. Используется для описания теплового состояния среды. \\
\hline
\texttt{ht.k\_iso} & Изотропная теплопроводность материала. Характеризует способность материала проводить тепло. \\
\hline
\texttt{ht.rho} & Плотность, используемая в тепловом модуле. Может использоваться при расчёте теплоёмкости на единицу объёма. \\
\hline
\texttt{ht.Cp} & Удельная теплоёмкость материала. Определяет количество энергии, необходимое для изменения температуры. \\
\hline
\texttt{u}, \texttt{v}, \texttt{w} & Компоненты перемещения узла по координатным осям. \\
\hline
\texttt{ut}, \texttt{vt}, \texttt{wt} & Компоненты скорости перемещения или временные производные перемещений, если они экспортированы из COMSOL. \\
\hline
\texttt{solid.E} & Модуль Юнга. Описывает жёсткость материала при упругой деформации. \\
\hline
\texttt{solid.nu} & Коэффициент Пуассона. Описывает связь продольной и поперечной деформации. \\
\hline
\texttt{solid.rho} & Плотность материала в механическом модуле. Может совпадать или отличаться от плотности, используемой в тепловой постановке. \\
\hline
\texttt{te1.alpha\_iso} & Изотропный коэффициент температурного расширения. Связывает изменение температуры с термической деформацией. \\
\hline
\texttt{solid.sx}, \texttt{solid.sy}, \texttt{solid.sz} & Нормальные компоненты тензора напряжений. \\
\hline
\texttt{solid.sxy}, \texttt{solid.syz}, \texttt{solid.sxz} & Касательные компоненты тензора напряжений. \\
\hline
\texttt{solid.eX}, \texttt{solid.eY}, \texttt{solid.eZ} & Компоненты деформации, экспортированные из механической части расчёта. \\
\hline
\end{tabular}
\end{table}

При подготовке данных важно привести все таблицы к единому формату. Для каждого узла и каждого временного шага должна быть сформирована одна строка или одна запись в структуре данных, содержащая координаты, физические поля и свойства материала. Если данные экспортированы отдельными CSV-файлами, они объединяются по координатам, идентификатору узла и времени. На практике предпочтительно использовать явный идентификатор узла, полученный из сетки, так как объединение только по координатам может быть чувствительно к округлению чисел с плавающей точкой.

\section{Преобразование конечно-элементной сетки в граф}

Пусть конечно-элементная сетка обозначается как $\mathcal{M}$. Для использования в MeshGraphNet она преобразуется в граф:
\begin{equation}
    G=(V,E),
\end{equation}
где $V$ --- множество узлов графа, а $E$ --- множество рёбер. Каждый узел $i \in V$ соответствует узлу конечно-элементной сетки. Рёбра строятся на основе связности конечных элементов.

Если два узла входят в один конечный элемент, между ними создаётся ребро. Для трёхмерной тетраэдральной сетки один конечный элемент обычно содержит четыре узла. Пусть тетраэдр задан набором узлов:
\begin{equation}
    \tau = (a,b,c,d).
\end{equation}
Тогда для неориентированного графа из этого тетраэдра получаются следующие пары соседей:
\begin{equation}
    \{(a,b),(a,c),(a,d),(b,c),(b,d),(c,d)\}.
\end{equation}
То есть каждый тетраэдр превращается в полный локальный подграф между своими четырьмя вершинами. Если граф используется в ориентированном виде, то для каждой пары добавляются два направления:
\begin{equation}
    (i,j) \in E \Rightarrow (j,i) \in E.
\end{equation}
Для одного тетраэдра это даёт двенадцать ориентированных рёбер:
\begin{equation}
\begin{split}
    E_{\tau}^{\mathrm{dir}}=\{&(a,b),(b,a),(a,c),(c,a),(a,d),(d,a),\\
    &(b,c),(c,b),(b,d),(d,b),(c,d),(d,c)\}.
\end{split}
\end{equation}

Ориентированное представление удобно для реализации message passing в PyTorch Geometric. В такой реализации граф обычно хранится через массив \texttt{edge\_index} размера $2 \times |E|$:
\begin{equation}
    \texttt{edge\_index}=\begin{bmatrix}
    i_1 & i_2 & \dots & i_{|E|} \\
    j_1 & j_2 & \dots & j_{|E|}
    \end{bmatrix},
\end{equation}
где каждый столбец задаёт одно ориентированное ребро $(i_k,j_k)$. Первая строка содержит индексы исходных узлов, вторая строка --- индексы целевых узлов.

При построении графа необходимо удалить дубликаты рёбер, потому что одна и та же пара узлов может входить в несколько соседних конечных элементов. После удаления дубликатов формируется итоговый список рёбер. Также важно сохранить соответствие между индексами узлов в файле сетки и строками в таблицах физических полей. Это соответствие используется для формирования матрицы признаков узлов.

\section{Узловые признаки}

Для каждого узла $i$ в момент времени $t$ формируется вектор входных признаков. В данной работе он включает координаты узла, тепловые величины, механические величины, свойства материала, напряжения и деформации:
\begin{equation}
\begin{split}
    h_i^t = [&x_i, y_i, z_i,
    T_i^t,
    u_i^t, v_i^t, w_i^t,
    ut_i^t, vt_i^t, wt_i^t,\\
    &E_i, \nu_i, \rho_i, \alpha_i, k_i, C_{p,i},
    \sigma_{x,i}^t, \sigma_{y,i}^t, \sigma_{z,i}^t,
    \sigma_{xy,i}^t, \sigma_{yz,i}^t, \sigma_{xz,i}^t,\\
    &\varepsilon_{x,i}^t, \varepsilon_{y,i}^t, \varepsilon_{z,i}^t].
\end{split}
\end{equation}

Вектор $h_i^t$ является строкой матрицы узловых признаков:
\begin{equation}
    H^t \in \mathbb{R}^{N \times F_h},
\end{equation}
где $N$ --- число узлов графа, $F_h$ --- число входных признаков на один узел.

Узловые признаки можно разделить на несколько групп.

\subsection{Координатные признаки}

Координатные признаки задают положение узла в трёхмерной области:
\begin{equation}
    p_i = [x_i,y_i,z_i].
\end{equation}
Они необходимы модели для понимания геометрического расположения узлов. Даже если граф уже содержит информацию о связности, координаты позволяют учитывать реальные расстояния, направление распространения теплового и механического воздействия, положение узла относительно нагретого стержня и границ области.

\subsection{Тепловые признаки}

К тепловым признакам относятся температура и тепловые параметры материала:
\begin{equation}
    h_{i,\mathrm{thermal}}^t=[T_i^t,k_i,C_{p,i},\rho_{i,\mathrm{thermal}}].
\end{equation}
Температура $T_i^t$ является основной динамической величиной тепловой части задачи. Теплопроводность $k_i$ влияет на скорость распространения тепла, теплоёмкость $C_{p,i}$ определяет тепловую инерцию материала, а плотность участвует в уравнении теплопереноса через объёмную теплоёмкость.

\subsection{Механические признаки}

Механические признаки описывают перемещение и скорость движения узлов:
\begin{equation}
    h_{i,\mathrm{mech}}^t=[u_i^t,v_i^t,w_i^t,ut_i^t,vt_i^t,wt_i^t].
\end{equation}
Компоненты $u_i^t$, $v_i^t$, $w_i^t$ задают смещение узла по осям координат. Компоненты $ut_i^t$, $vt_i^t$, $wt_i^t$ описывают скорость изменения перемещений, если они экспортированы из COMSOL. Для волновых и термоупругих процессов скорости важны, так как они несут информацию о динамическом состоянии системы и направлении распространения механического возмущения.

\subsection{Материальные признаки}

Материальные признаки описывают локальные свойства среды:
\begin{equation}
    h_{i,\mathrm{mat}}=[E_i,\nu_i,\rho_i,\alpha_i,k_i,C_{p,i}].
\end{equation}
Здесь $E_i$ --- модуль Юнга, $\nu_i$ --- коэффициент Пуассона, $\rho_i$ --- плотность, $\alpha_i$ --- коэффициент температурного расширения, $k_i$ --- теплопроводность, $C_{p,i}$ --- удельная теплоёмкость. Эти признаки особенно важны, если модель должна обучаться на нескольких симуляциях с различными породами или различными параметрами металлического стержня. В этом случае MeshGraphNet получает возможность учитывать физические различия материалов через входные данные, а не только через обучаемые веса.

\subsection{Признаки напряжений}

Напряжённое состояние узла описывается компонентами тензора напряжений:
\begin{equation}
    h_{i,\sigma}^t=[\sigma_{x,i}^t,\sigma_{y,i}^t,\sigma_{z,i}^t,\sigma_{xy,i}^t,\sigma_{yz,i}^t,\sigma_{xz,i}^t].
\end{equation}
В обозначениях COMSOL этим компонентам соответствуют поля \texttt{solid.sx}, \texttt{solid.sy}, \texttt{solid.sz}, \texttt{solid.sxy}, \texttt{solid.syz}, \texttt{solid.sxz}. Эти величины позволяют модели учитывать механический отклик среды на температурное воздействие.

\subsection{Признаки деформаций}

Деформации описывают относительное изменение формы материала:
\begin{equation}
    h_{i,\varepsilon}^t=[\varepsilon_{x,i}^t,\varepsilon_{y,i}^t,\varepsilon_{z,i}^t].
\end{equation}
В экспортированных данных они могут соответствовать столбцам \texttt{solid.eX}, \texttt{solid.eY}, \texttt{solid.eZ}. Деформации являются связующим звеном между перемещениями и напряжениями, поэтому их использование в признаках помогает модели точнее описывать механическую часть термоупругого процесса.

Сводное описание групп узловых признаков приведено в таблице~\ref{tab:node_features}.

\begin{table}[h]
\centering
\caption{Группы узловых признаков MeshGraphNet}
\label{tab:node_features}
\begin{tabular}{|p{0.25\textwidth}|p{0.32\textwidth}|p{0.33\textwidth}|}
\hline
\textbf{Группа} & \textbf{Поля} & \textbf{Назначение} \\
\hline
Координаты & $x$, $y$, $z$ & Геометрическое положение узла в расчётной области. \\
\hline
Температура & $T$ & Текущее тепловое состояние узла. \\
\hline
Перемещения и скорости & $u$, $v$, $w$, $ut$, $vt$, $wt$ & Динамическое механическое состояние среды. \\
\hline
Материалы & $E$, $\nu$, $\rho$, $\alpha$, $k$, $C_p$ & Локальные физические свойства породы или металлического стержня. \\
\hline
Напряжения & $\sigma_x$, $\sigma_y$, $\sigma_z$, $\sigma_{xy}$, $\sigma_{yz}$, $\sigma_{xz}$ & Описание напряжённого состояния материала. \\
\hline
Деформации & $\varepsilon_x$, $\varepsilon_y$, $\varepsilon_z$ & Описание локальной деформации среды. \\
\hline
\end{tabular}
\end{table}

\section{Рёберные признаки}

Рёбра графа описывают локальную связь между соседними узлами конечно-элементной сетки. Для каждого ориентированного ребра $(i,j)$ формируется вектор рёберных признаков:
\begin{equation}
    r_{ij}=[x_j-x_i,\;y_j-y_i,\;z_j-z_i,\;l_{ij}],
\end{equation}
где $l_{ij}$ --- длина ребра:
\begin{equation}
    l_{ij}=\sqrt{(x_j-x_i)^2+(y_j-y_i)^2+(z_j-z_i)^2}.
\end{equation}

Относительные координаты $(x_j-x_i,y_j-y_i,z_j-z_i)$ задают направление от узла $i$ к узлу $j$. Длина ребра $l_{ij}$ задаёт расстояние между узлами. Эти признаки важны по физическим причинам. Тепловое и механическое влияние зависит от расстояния между точками: при прочих равных условиях соседние узлы, расположенные ближе друг к другу, сильнее связаны в локальной дискретизации. Кроме того, направление ребра помогает модели учитывать пространственную ориентацию распространения возмущений.

Матрица рёберных признаков обозначается как:
\begin{equation}
    R \in \mathbb{R}^{|E| \times F_r},
\end{equation}
где $|E|$ --- число ориентированных рёбер, $F_r$ --- число признаков ребра. В базовом варианте $F_r=4$. При необходимости рёберный вектор может быть расширен дополнительными признаками, например типом контакта, принадлежностью ребра границе или принадлежностью ребра определённой материальной области.

\section{Временная структура данных}

Так как в COMSOL при экспорте может использоваться режим \texttt{Time selection: All}, физические поля сохраняются для нескольких временных шагов. Поэтому датасет MeshGraphNet имеет не только пространственную, но и временную структуру.

Состояние графа в момент времени $t$ можно записать как:
\begin{equation}
    G^t=(V,E,H^t,R),
\end{equation}
где $V$ и $E$ задают неизменную топологию сетки, $H^t$ --- матрица узловых признаков в момент времени $t$, а $R$ --- матрица рёберных признаков. В большинстве симуляций сетка не меняется во времени, поэтому $V$, $E$ и $R$ являются статическими, а изменяются только физические поля в $H^t$.

Цель модели --- спрогнозировать состояние в следующий момент времени:
\begin{equation}
    Y^{t+1}=F_{\theta}(G^t),
\end{equation}
где $Y^{t+1}$ --- матрица целевых физических величин в момент времени $t+1$.

Возможны два способа постановки задачи. Первый способ --- прогнозировать абсолютные значения физических величин:
\begin{equation}
    \hat{Y}^{t+1}=F_{\theta}(G^t).
\end{equation}
Второй способ --- прогнозировать приращения:
\begin{equation}
    \Delta Y^t=Y^{t+1}-Y^t,
\end{equation}
\begin{equation}
    \widehat{\Delta Y}^t=F_{\theta}(G^t),
\end{equation}
\begin{equation}
    \hat{Y}^{t+1}=Y^t+\widehat{\Delta Y}^t.
\end{equation}

Прогнозирование приращений часто удобно для динамических физических задач. Во-первых, приращения обычно имеют меньший масштаб, чем абсолютные значения. Во-вторых, модель обучается предсказывать локальное изменение состояния, что ближе к идее численного интегрирования по времени. В-третьих, такой подход может стабилизировать rollout-прогнозирование, так как модель явно обновляет текущее состояние, а не генерирует новое состояние с нуля.

\section{Целевые переменные}

Целевой вектор задаёт, какие физические величины MeshGraphNet должна прогнозировать. В полной постановке для узла $i$ в момент времени $t+1$ можно использовать следующий вектор:
\begin{equation}
\begin{split}
    y_i^{t+1}=[&T_i^{t+1},
    u_i^{t+1},v_i^{t+1},w_i^{t+1},
    \sigma_{x,i}^{t+1},\sigma_{y,i}^{t+1},\sigma_{z,i}^{t+1},\\
    &\sigma_{xy,i}^{t+1},\sigma_{yz,i}^{t+1},\sigma_{xz,i}^{t+1}].
\end{split}
\end{equation}

Тогда матрица целевых значений имеет вид:
\begin{equation}
    Y^{t+1}\in\mathbb{R}^{N\times F_y},
\end{equation}
где $F_y$ --- число прогнозируемых физических каналов.

В зависимости от цели эксперимента набор целевых переменных может быть уменьшен. Например, возможны следующие варианты:
\begin{itemize}
    \item прогноз только температуры: $y_i^{t+1}=[T_i^{t+1}]$;
    \item прогноз температуры и перемещений: $y_i^{t+1}=[T_i^{t+1},u_i^{t+1},v_i^{t+1},w_i^{t+1}]$;
    \item прогноз температуры, перемещений и напряжений: $y_i^{t+1}=[T_i^{t+1},u_i^{t+1},v_i^{t+1},w_i^{t+1},\sigma_i^{t+1}]$.
\end{itemize}

Выбор целевых переменных влияет на сложность задачи. Прогноз только температуры является более простой постановкой, так как модель учится описывать тепловое поле. Прогноз температуры и перемещений уже требует учитывать термомеханическую связь. Прогноз напряжений является более сложным, потому что напряжения зависят от деформаций, свойств материала, температурного расширения и граничных условий.

\section{Архитектура MeshGraphNet}

MeshGraphNet строится по схеме \textit{Encode--Process--Decode}. Такая архитектура состоит из трёх основных частей:
\begin{enumerate}
    \item \textit{Encoder} --- переводит исходные узловые и рёберные признаки в латентное пространство;
    \item \textit{Processor} --- выполняет несколько итераций передачи сообщений по графу;
    \item \textit{Decoder} --- преобразует итоговые латентные признаки узлов в прогнозируемые физические величины.
\end{enumerate}

\subsection{Encoder}

На вход encoder подаются матрица узловых признаков $H^t$ и матрица рёберных признаков $R$. Для каждого узла и каждого ребра применяются отдельные многослойные перцептроны:
\begin{equation}
    z_i^0=f_v^{\mathrm{enc}}(h_i^t),
\end{equation}
\begin{equation}
    q_{ij}^0=f_e^{\mathrm{enc}}(r_{ij}),
\end{equation}
где $z_i^0\in\mathbb{R}^{H}$ --- начальное латентное представление узла, $q_{ij}^0\in\mathbb{R}^{H}$ --- начальное латентное представление ребра, $H$ --- размерность скрытого пространства. Функции $f_v^{\mathrm{enc}}$ и $f_e^{\mathrm{enc}}$ обычно реализуются как MLP с нелинейными функциями активации и нормализацией.

Назначение encoder состоит в том, чтобы привести признаки разной природы к единому скрытому представлению. Например, температура измеряется в кельвинах, перемещения --- в метрах, напряжения --- в паскалях, а координаты --- в единицах длины. Encoder формирует из этих величин латентные векторы, с которыми далее работает processor.

\subsection{Processor}

Processor является основной частью MeshGraphNet. Он выполняет $K$ итераций message passing. На каждой итерации обновляются признаки рёбер, затем сообщения агрегируются в узлах, после чего обновляются признаки узлов.

Для итерации $k$ обновление ребра можно записать как:
\begin{equation}
    q_{ij}^{k+1}=\phi_e\left(q_{ij}^{k},z_i^k,z_j^k\right),
\end{equation}
где $z_i^k$ и $z_j^k$ --- латентные признаки узлов на концах ребра, $q_{ij}^{k}$ --- текущее латентное представление ребра, $\phi_e$ --- обучаемая функция обновления ребра.

После обновления рёбер для каждого узла выполняется агрегация входящих сообщений:
\begin{equation}
    m_i^{k+1}=\sum_{j\in\mathcal{N}(i)}q_{ji}^{k+1},
\end{equation}
где $\mathcal{N}(i)$ --- множество соседей, из которых в узел $i$ приходят сообщения. Суммирование обеспечивает инвариантность к порядку соседей: результат не зависит от того, в каком порядке рёбра перечислены в массиве \texttt{edge\_index}.

Затем обновляется латентное представление узла:
\begin{equation}
    z_i^{k+1}=\phi_v\left(z_i^k,m_i^{k+1}\right),
\end{equation}
где $\phi_v$ --- обучаемая функция обновления узла.

На практике часто используется остаточная форма обновления:
\begin{equation}
    q_{ij}^{k+1}=q_{ij}^{k}+\phi_e\left(q_{ij}^{k},z_i^k,z_j^k\right),
\end{equation}
\begin{equation}
    z_i^{k+1}=z_i^k+\phi_v\left(z_i^k,m_i^{k+1}\right).
\end{equation}
Остаточные связи улучшают устойчивость обучения, особенно при большом количестве итераций message passing. Также могут использоваться LayerNorm, dropout и другие стандартные приёмы стабилизации обучения нейронных сетей.

\subsection{Decoder}

После $K$ итераций processor каждый узел имеет итоговое латентное представление $z_i^K$. Decoder преобразует его в прогноз физического состояния:
\begin{equation}
    \hat{y}_i^{t+1}=f_{\mathrm{dec}}(z_i^K).
\end{equation}
Если модель обучается прогнозировать приращение, decoder выдаёт не само значение $y_i^{t+1}$, а изменение состояния:
\begin{equation}
    \widehat{\Delta y}_i^t=f_{\mathrm{dec}}(z_i^K),
\end{equation}
после чего прогноз следующего шага получается как:
\begin{equation}
    \hat{y}_i^{t+1}=y_i^t+\widehat{\Delta y}_i^t.
\end{equation}

\section{Message passing и физический смысл передачи сообщений}

Message passing в MeshGraphNet можно интерпретировать как нейросетевый аналог локального обмена информацией между соседними узлами конечно-элементной сетки. В физической задаче температура, перемещения, деформации и напряжения в одном узле зависят не только от состояния самого узла, но и от состояния окружающих узлов. Например, тепловой поток определяется разностью температур и свойствами материала между соседними точками, а механическое напряжение связано с локальными деформациями и взаимодействием соседних областей среды.

На одном шаге message passing узел получает информацию только от непосредственных соседей. После двух шагов он получает информацию уже от соседей второго порядка, после трёх --- от более широкой окрестности. В общем случае после $K$ итераций узел может учитывать информацию из $K$-hop окрестности:
\begin{equation}
    z_i^K = \mathrm{MP}^{K}\left(z_i^0,\{z_j^0: j\in\mathcal{N}_K(i)\}\right),
\end{equation}
где $\mathcal{N}_K(i)$ --- множество узлов, достижимых из узла $i$ не более чем за $K$ рёбер.

Физически это похоже на распространение локального влияния по сетке. Чем больше число message passing слоёв, тем более широкую область может учитывать модель при прогнозе состояния каждого узла. Однако слишком большое число слоёв может привести к росту вычислительной стоимости, переусреднению латентных признаков и усложнению обучения. Поэтому параметр $K$ подбирается экспериментально с учётом размера сетки, временного шага и требуемой точности прогноза.

Для термоупругой задачи такая локальная передача сообщений особенно важна. Температурное поле влияет на тепловое расширение, тепловое расширение влияет на деформации, деформации связаны с напряжениями, а напряжения и перемещения распространяются по среде в виде механического отклика. MeshGraphNet не задаёт эти уравнения явно, но учится воспроизводить их дискретную динамику на основе данных COMSOL.

\section{Нормализация данных}

Физические величины в задаче имеют разные масштабы. Температура может измеряться в сотнях кельвинов, координаты и перемещения --- в метрах, напряжения --- в паскалях, модуль Юнга --- в гигапаскалях, коэффициент температурного расширения --- в малых значениях порядка $10^{-6}$. Если подавать такие величины в нейронную сеть без нормализации, признаки с большими численными значениями могут доминировать над остальными, а обучение становится менее устойчивым.

Поэтому для каждого канала признаков применяется стандартизация:
\begin{equation}
    \tilde{x}=\frac{x-\mu}{\sigma},
\end{equation}
где $\mu$ --- среднее значение признака, $\sigma$ --- стандартное отклонение признака. Статистики $\mu$ и $\sigma$ вычисляются только по обучающей выборке. Это важно, потому что использование validation или test данных при расчёте статистик приводит к утечке информации из тестовой части в обучение.

Для разных групп признаков статистики считаются отдельно:
\begin{itemize}
    \item отдельно для координат $x$, $y$, $z$;
    \item отдельно для температуры $T$;
    \item отдельно для перемещений $u$, $v$, $w$;
    \item отдельно для скоростей $ut$, $vt$, $wt$;
    \item отдельно для напряжений $\sigma_x$, $\sigma_y$, $\sigma_z$, $\sigma_{xy}$, $\sigma_{yz}$, $\sigma_{xz}$;
    \item отдельно для деформаций $\varepsilon_x$, $\varepsilon_y$, $\varepsilon_z$;
    \item отдельно для материальных параметров $E$, $\nu$, $\rho$, $\alpha$, $k$, $C_p$;
    \item отдельно для целевых переменных или их приращений.
\end{itemize}

Если модель прогнозирует приращения, то нормализация выходов выполняется для величин $\Delta Y^t$, а не только для абсолютных значений $Y^t$. Это позволяет привести масштаб прогнозируемых изменений к удобному диапазону и улучшить сходимость обучения.

При обратном переходе к физическим единицам выполняется денормализация:
\begin{equation}
    x=\tilde{x}\sigma+\mu.
\end{equation}
Она используется при сохранении предсказаний, построении графиков, расчёте метрик в физических единицах и сравнении прогноза с результатами COMSOL.

\section{Разделение датасета}

Датасет MeshGraphNet должен разделяться на обучающую, валидационную и тестовую части. При этом важно разделять данные по симуляциям, а не случайно по отдельным строкам таблиц. Причина состоит в том, что строки одной симуляции сильно связаны между собой: они относятся к одной геометрии, одной сетке, одной последовательности временных шагов и одному набору физических параметров.

Если случайно перемешать строки или временные шаги одной симуляции между train и test, модель может фактически увидеть в обучении почти те же состояния, которые затем используются для проверки. В этом случае оценка качества будет завышенной и не будет отражать способность модели обобщать результат на новые симуляции.

Формально полный набор симуляций можно обозначить как:
\begin{equation}
    \mathcal{D}=\{S_1,S_2,\dots,S_n\},
\end{equation}
где $S_k$ --- одна симуляция, содержащая сетку, параметры материалов и последовательность временных графов. Разделение выполняется так:
\begin{equation}
    \mathcal{D}=\mathcal{D}_{\mathrm{train}}\cup\mathcal{D}_{\mathrm{val}}\cup\mathcal{D}_{\mathrm{test}},
\end{equation}
\begin{equation}
    \mathcal{D}_{\mathrm{train}}\cap\mathcal{D}_{\mathrm{val}}=\varnothing,
\end{equation}
\begin{equation}
    \mathcal{D}_{\mathrm{train}}\cap\mathcal{D}_{\mathrm{test}}=\varnothing,
\end{equation}
\begin{equation}
    \mathcal{D}_{\mathrm{val}}\cap\mathcal{D}_{\mathrm{test}}=\varnothing.
\end{equation}

Если доступна только одна длинная симуляция, то разделение становится менее надёжным. В таком случае можно разделять данные по непрерывным временным интервалам, например первые временные шаги использовать для обучения, следующие --- для валидации, последние --- для теста. Однако такой подход всё равно проверяет в основном интерполяцию по времени внутри одной физической постановки, а не полноценное обобщение на новые параметры материала или новую температуру стержня. Поэтому для более корректного обучения желательно иметь несколько симуляций с различными параметрами породы, начальной температурой стержня или другими физическими условиями.

\section{Функция потерь}

Базовой функцией потерь является среднеквадратичная ошибка по узлам и прогнозируемым физическим каналам:
\begin{equation}
    \mathcal{L}_{\mathrm{MSE}}=\frac{1}{N}\sum_{i=1}^{N}\left\|\hat{y}_i-y_i\right\|_2^2.
\end{equation}
Если используется батч из нескольких графов, ошибка усредняется по всем узлам всех графов батча.

Для термоупругой задачи разные группы целевых переменных имеют разные физические единицы и разные масштабы. Поэтому может использоваться взвешенная функция потерь:
\begin{equation}
    \mathcal{L}=w_T\mathcal{L}_T+w_u\mathcal{L}_u+w_{\sigma}\mathcal{L}_{\sigma},
\end{equation}
где $\mathcal{L}_T$ --- ошибка температуры, $\mathcal{L}_u$ --- ошибка перемещений, $\mathcal{L}_{\sigma}$ --- ошибка напряжений, а $w_T$, $w_u$, $w_{\sigma}$ --- веса соответствующих частей loss.

Ошибка температуры может быть задана как:
\begin{equation}
    \mathcal{L}_T=\frac{1}{N}\sum_{i=1}^{N}\left(\hat{T}_i-T_i\right)^2.
\end{equation}

Ошибка перемещений может быть задана как средняя ошибка по трём компонентам:
\begin{equation}
    \mathcal{L}_u=\frac{1}{N}\sum_{i=1}^{N}\left[\left(\hat{u}_i-u_i\right)^2+\left(\hat{v}_i-v_i\right)^2+\left(\hat{w}_i-w_i\right)^2\right].
\end{equation}

Ошибка напряжений может быть задана как средняя ошибка по шести компонентам тензора напряжений:
\begin{equation}
\begin{split}
    \mathcal{L}_{\sigma}=\frac{1}{N}\sum_{i=1}^{N}[&
    (\hat{\sigma}_{x,i}-\sigma_{x,i})^2+
    (\hat{\sigma}_{y,i}-\sigma_{y,i})^2+
    (\hat{\sigma}_{z,i}-\sigma_{z,i})^2+\\
    &(\hat{\sigma}_{xy,i}-\sigma_{xy,i})^2+
    (\hat{\sigma}_{yz,i}-\sigma_{yz,i})^2+
    (\hat{\sigma}_{xz,i}-\sigma_{xz,i})^2].
\end{split}
\end{equation}

Веса $w_T$, $w_u$ и $w_{\sigma}$ нужны для балансировки вклада разных физических величин. Например, если напряжения имеют гораздо больший численный масштаб, чем перемещения, то без нормализации и весов модель может оптимизировать главным образом напряжения и хуже обучаться температуре или перемещениям. При корректной нормализации необходимость в ручном подборе весов уменьшается, но взвешенная loss всё равно может быть полезна, если некоторые физические поля имеют больший приоритет в исследовании.

\section{Обучение модели}

Процесс обучения MeshGraphNet состоит из последовательности стандартных этапов машинного обучения, адаптированных к графовым данным. Один обучающий пример соответствует паре графов или паре состояний одной симуляции:
\begin{equation}
    (G^t,Y^{t+1}).
\end{equation}
Если модель прогнозирует приращения, целевым значением является:
\begin{equation}
    \Delta Y^t=Y^{t+1}-Y^t.
\end{equation}

Общий алгоритм обучения можно описать следующим образом:
\begin{enumerate}
    \item загрузить конечно-элементную сетку из \texttt{sandstone.mphtxt};
    \item построить список узлов, список рёбер и массив \texttt{edge\_index};
    \item загрузить CSV-файлы с физическими полями;
    \item объединить данные по узлам и временным шагам;
    \item сформировать матрицы узловых признаков $H^t$, рёберных признаков $R$ и целевых значений $Y^{t+1}$;
    \item применить нормализацию признаков и целевых переменных;
    \item передать граф $G^t$ в encoder;
    \item выполнить $K$ итераций message passing в processor;
    \item получить прогноз через decoder;
    \item вычислить функцию потерь между прогнозом и данными COMSOL;
    \item выполнить обратное распространение ошибки;
    \item обновить веса модели с помощью оптимизатора.
\end{enumerate}

При реализации на PyTorch Geometric один граф может храниться в структуре данных, содержащей:
\begin{itemize}
    \item \texttt{x} --- матрицу узловых признаков;
    \item \texttt{edge\_index} --- индексы рёбер;
    \item \texttt{edge\_attr} --- матрицу рёберных признаков;
    \item \texttt{y} --- целевые значения;
    \item \texttt{pos} --- координаты узлов;
    \item \texttt{time\_idx} --- индекс временного шага;
    \item \texttt{simulation\_id} --- идентификатор симуляции.
\end{itemize}

Для ускорения обучения используется графический процессор. Если доступна видеокарта с поддержкой CUDA, модель, батчи графов и тензоры переносятся на GPU. Это особенно важно для больших конечно-элементных сеток, так как число узлов и рёбер может быть значительным.

Основные гиперпараметры обучения приведены в таблице~\ref{tab:training_params}.

\begin{table}[h]
\centering
\caption{Основные гиперпараметры обучения MeshGraphNet}
\label{tab:training_params}
\begin{tabular}{|p{0.28\textwidth}|p{0.62\textwidth}|}
\hline
\textbf{Параметр} & \textbf{Назначение} \\
\hline
\texttt{batch\_size} & Число графов или временных переходов, обрабатываемых за один шаг оптимизации. \\
\hline
\texttt{hidden\_size} & Размерность латентных векторов узлов и рёбер. \\
\hline
\texttt{message\_passing\_steps} & Число итераций processor, определяющее размер учитываемой окрестности узла. \\
\hline
\texttt{learning\_rate} & Скорость обучения оптимизатора. \\
\hline
\texttt{epochs} & Количество полных проходов по обучающей выборке. \\
\hline
\texttt{weight\_decay} & Регуляризация весов модели. \\
\hline
\texttt{target\_mode} & Режим прогноза: абсолютные значения или приращения. \\
\hline
\end{tabular}
\end{table}

Во время обучения контролируется ошибка на validation-наборе. Лучшая версия модели сохраняется в файл checkpoint. Затем этот checkpoint используется для инференса и rollout-прогнозирования.

\section{Rollout-прогнозирование}

Одношаговый прогноз соответствует применению модели к одному состоянию графа:
\begin{equation}
    \hat{Y}^{t+1}=F_{\theta}(G^t).
\end{equation}
Такой режим используется при обучении и первичной оценке качества модели. Однако для динамической физической задачи важен не только один шаг, но и способность модели последовательно прогнозировать развитие процесса во времени.

Многошаговый прогноз, или rollout, выполняется рекурсивно. Сначала модель получает истинное состояние $G^t$ и предсказывает $\hat{Y}^{t+1}$. Затем это предсказанное состояние используется как вход для следующего шага:
\begin{equation}
    \hat{Y}^{t+1}=F_{\theta}(G^t),
\end{equation}
\begin{equation}
    \hat{Y}^{t+2}=F_{\theta}(\hat{G}^{t+1}),
\end{equation}
\begin{equation}
    \hat{Y}^{t+3}=F_{\theta}(\hat{G}^{t+2}).
\end{equation}
В общем виде:
\begin{equation}
    \hat{Y}^{t+s}=F_{\theta}(\hat{G}^{t+s-1}), \qquad s=1,2,\dots,S.
\end{equation}

Если используется прогноз приращений, обновление состояния записывается как:
\begin{equation}
    \hat{Y}^{t+s}=\hat{Y}^{t+s-1}+F_{\theta}(\hat{G}^{t+s-1}).
\end{equation}

Rollout важен, потому что он показывает способность модели воспроизводить динамику термоупругого процесса на длительном временном интервале. При этом ошибка может накапливаться: небольшая ошибка на первом шаге изменяет вход для второго шага, затем ошибка второго шага влияет на третий и так далее. Поэтому качество rollout-прогноза обычно ниже качества одношагового прогноза, но именно rollout лучше отражает практическую применимость модели для симуляции распространения волн и изменения температурно-напряжённого состояния среды.

\section{Метрики качества}

Качество MeshGraphNet оценивается относительно расчётных данных COMSOL. Цель оценки --- определить, насколько точно модель воспроизводит температурные, механические и напряжённые поля, полученные численным решателем.

Средняя абсолютная ошибка задаётся как:
\begin{equation}
    \mathrm{MAE}=\frac{1}{N}\sum_{i=1}^{N}|\hat{y}_i-y_i|.
\end{equation}

Среднеквадратичная ошибка задаётся как:
\begin{equation}
    \mathrm{MSE}=\frac{1}{N}\sum_{i=1}^{N}(\hat{y}_i-y_i)^2.
\end{equation}

Корень из среднеквадратичной ошибки задаётся как:
\begin{equation}
    \mathrm{RMSE}=\sqrt{\frac{1}{N}\sum_{i=1}^{N}(\hat{y}_i-y_i)^2}.
\end{equation}

Относительная ошибка может быть записана как:
\begin{equation}
    \mathrm{RelError}=\frac{\left\|\hat{Y}-Y\right\|_2}{\left\|Y\right\|_2+\epsilon},
\end{equation}
где $\epsilon$ --- малое число для предотвращения деления на ноль.

Для отдельных групп физических величин используются специализированные метрики. Ошибка температуры:
\begin{equation}
    \mathrm{RMSE}_T=\sqrt{\frac{1}{N}\sum_{i=1}^{N}(\hat{T}_i-T_i)^2}.
\end{equation}

Ошибка перемещений может считаться по норме вектора перемещения:
\begin{equation}
    \mathrm{RMSE}_u=\sqrt{\frac{1}{N}\sum_{i=1}^{N}\left\|\hat{\mathbf{u}}_i-\mathbf{u}_i\right\|_2^2},
\end{equation}
где
\begin{equation}
    \mathbf{u}_i=[u_i,v_i,w_i].
\end{equation}

Ошибка напряжений может считаться по компонентам тензора напряжений:
\begin{equation}
    \mathrm{RMSE}_{\sigma}=\sqrt{\frac{1}{N}\sum_{i=1}^{N}\left\|\hat{\boldsymbol{\sigma}}_i-\boldsymbol{\sigma}_i\right\|_2^2},
\end{equation}
где
\begin{equation}
    \boldsymbol{\sigma}_i=[\sigma_{x,i},\sigma_{y,i},\sigma_{z,i},\sigma_{xy,i},\sigma_{yz,i},\sigma_{xz,i}].
\end{equation}

Также полезно оценивать ошибку максимальных значений, так как в инженерных задачах часто важны пиковые температуры и пиковые напряжения:
\begin{equation}
    E_{T,\max}=\left|\max_i \hat{T}_i-\max_i T_i\right|,
\end{equation}
\begin{equation}
    E_{\sigma,\max}=\left|\max_i \left\|\hat{\boldsymbol{\sigma}}_i\right\|_2-\max_i \left\|\boldsymbol{\sigma}_i\right\|_2\right|.
\end{equation}

Эти метрики используются не для сравнения MeshGraphNet с другими алгоритмами, а для проверки качества нейросетевой аппроксимации относительно COMSOL-данных.

\section{Итог по MeshGraphNet-части}

MeshGraphNet-часть дипломной работы предназначена для построения графовой нейросетевой модели, которая обучается на результатах конечно-элементных симуляций COMSOL. Файл \texttt{sandstone.mphtxt} используется как источник конечно-элементной сетки, из которой формируются узлы, рёбра и рёберные признаки графа. CSV-файлы \texttt{data\_temperature.csv}, \texttt{data\_displacement.csv}, \texttt{data\_materials.csv}, \texttt{data\_strain.csv}, \texttt{data\_stress\_1.csv}, \texttt{data\_stress\_2.csv} и \texttt{data\_stress\_3.csv} используются как источник физических полей и параметров материалов.

В рамках этой части данные COMSOL преобразуются в последовательность графов $G^t=(V,E,H^t,R)$, где узлы соответствуют узлам конечно-элементной сетки, рёбра описывают локальную связность сетки, узловые признаки содержат координаты, температуру, перемещения, скорости, свойства материалов, напряжения и деформации, а рёберные признаки содержат относительные координаты и расстояния между соседними узлами.

MeshGraphNet обучается прогнозировать следующее состояние физических полей или их приращение. Архитектура модели состоит из encoder, processor и decoder. Processor выполняет message passing, благодаря чему модель учитывает локальное распространение теплового и механического влияния по сетке. После обучения модель может использоваться для одношагового прогноза и rollout-прогнозирования динамики термоупругого процесса.

Таким образом, данная часть является самостоятельным машинно-обучающим модулем общей дипломной работы. Она связывает расчётные данные COMSOL с графовой нейронной сетью и позволяет исследовать возможность AI-направленного прогнозирования распространения термоупругих волн в геологических средах.
